\chapter{Project Introduction}

The rapid growth of online food delivery services has significantly transformed the way customers order and access food in modern lifestyles. With increasing dependence on digital platforms, users now expect fast, convenient, and personalized food ordering experiences.

\vspace{0.5mm}

Traditional food ordering systems primarily focus on basic order placement and delivery functionalities, often lacking intelligent interaction, personalized recommendations, and efficient management features. As competition in the food delivery industry continues to increase, there is a growing need for a smart and scalable system that can improve customer experience while simplifying restaurant and delivery operations.

\vspace{0.5mm}

The SwiftBite project explores these challenges through a web-based food-ordering prototype with a customer storefront, account flows, and demonstration portals for different user roles. It also provides chatbot assistance for common application questions. The prototype combines a React frontend with a backend API connected to MySQL; the database connection and restaurant API have been verified. Persistent ordering, payment processing, and production role authorization remain areas for further development.

\section{Project Objectives}

The primary objective of the SwiftBite project is to develop and evaluate a user-friendly food-ordering prototype with an account flow, a customer storefront, and illustrative interfaces for customer, restaurant, delivery, and administration workflows.

The project also includes chatbot assistance for application questions. A Gemini-generated recipe response was verified during testing; transient primary-model failures are retried and a secondary model is attempted before local rule-based guidance is used. The chatbot is not a personalized recommendation engine.

The project uses a React, TypeScript, and Vite frontend with a Node.js and Express API and a MySQL schema. The customer restaurant catalogue, cart, checkout, and tracking screens are currently prototype flows, rather than a complete real-time ordering and delivery service.

\section{Background of the Project}

The inspiration for SwiftBite arises from the demand for convenient online food-ordering experiences. The prototype explores this area through a customer-facing interface, account flows, and application guidance from a chatbot. It does not analyze customer preferences or generate personalized food recommendations.

The current implementation uses React and Vite for the frontend, an Express API for backend services, and a connected MySQL database with an initialized schema. The backend restaurant endpoint has been verified against MySQL. The backend Gemini integration returned a generated response during chatbot testing, with transient-error retry, a secondary model, and local rule-based guidance as fallback layers. The customer storefront still uses bundled catalogue data. Checkout and order status are demonstrations; the delivery screen provides a zoomable OpenStreetMap view centered on Bhubaneswar with a static marker. Persistent customer orders, payment processing, live rider tracking, personalized recommendations, and production access control remain future work.

\section{Operation Environment}

The SwiftBite application operates as a web-based prototype. The frontend is developed using \textbf{React}, \textbf{TypeScript}, and \textbf{Vite}, with \textbf{Tailwind CSS} and custom CSS for its responsive interface. The project also includes a web app manifest and service-worker registration.

The backend is developed using \textbf{Node.js}, \textbf{Express}, and \textbf{TypeScript}. It exposes REST API routes for authentication, restaurant and profile data, menu items, image uploads, orders, dashboards, and chatbot assistance. The backend is connected to \textbf{MySQL} through \texttt{mysql2}; its schema is initialized, and the restaurant API has been verified against the database. The frontend storefront currently uses bundled sample restaurant and menu data, and checkout does not persist orders or payment transactions.

Backend libraries include \textbf{CORS}, \textbf{Helmet}, \textbf{Morgan}, \textbf{Dotenv}, \textbf{mysql2}, and \textbf{Zod}. The chatbot uses the Google Generative AI SDK through a server-side key; local responses remain a fallback when Gemini is unavailable. Role-specific portal switching is marked as a demonstration feature, not production access control.